\chapter*{Závěr}
\addcontentsline{toc}{chapter}{Závěr}
\label{chap:zaver}

Cílem této práce bylo navrhnout a~implementovat mobilní aplikaci pro Android,
která umožní pokročilou správu Bitcoinových prostředků -- konkrétně podporu
multisignature transakcí ve schématu M-of-N, integraci s~hardwarovou peněženkou
Trezor a~manuální výběr UTXO -- a~zaplnit tak mezeru, kterou jsme v~analýze
existujících peněženek (viz~kapitolu~\ref{chap:related}) identifikovali na
mobilních platformách. Všechny tři stanovené cíle se podařilo naplnit;
ověření jejich funkčnosti popisujeme v~kapitole~\ref{chap:testovani}, kde jsme
mimo jiné provedli kompletní 2-of-3 multisig transakci od importu deskriptoru
přes postupné podepisování dvěma cosignery až po broadcast na Bitcoin testnet.

Výsledná aplikace se skládá z~šesti backendových mikroslužeb a~API Gateway
implementovaných v~jazyce Kotlin s~frameworkem Ktor~\cite{KtorFramework}
a~nativního Android frontendu postaveného na Jetpack Compose~\cite{JetpackCompose}.
Bezpečnostní model důsledně dodržuje princip, že server nikdy nedisponuje
privátními klíči ani jinými materiály, ze~kterých by je bylo možné odvodit --
veškeré podepisování probíhá výhradně na hardwarové peněžence, kterou
oslovujeme prostřednictvím deeplink protokolu Trezor Connect~\cite{TrezorConnect}
zprostředkovaného aplikací Trezor Suite Mobile~\cite{TrezorSuite}.
Pro distribuci částečně podepsaných transakcí mezi cosignery
využíváme standard PSBT~\cite{BIP174} a~konfiguraci multisig peněženky
přebíráme z~výstupních deskriptorů ve~formátu \texttt{wsh(sortedmulti(M,\ldots))}
podle BIP-380~\cite{BIP380} a~BIP-383~\cite{BIP383}. Tato volba zajišťuje
interoperabilitu s~existujícími desktopovými peněženkami -- typicky lze
vytvořit multisig konfiguraci ve Sparrow Wallet~\cite{SparrowWallet}, exportovat
deskriptor a~v~naší aplikaci jej importovat, čímž obě peněženky pracují nad
totožnou množinou adres.

Práce ukazuje, že kombinace pokročilých funkcí, kterou na desktopu poskytuje
Sparrow Wallet, je realizovatelná i~na mobilní platformě bez kompromisu
v~bezpečnostním modelu. Mobilní zařízení v~tomto uspořádání slouží jako
plnohodnotný koordinátor multisig transakcí (vytváří PSBT, sleduje stav
shromážděných podpisů, koordinuje broadcast), zatímco veškeré kryptografické
operace zůstávají izolovány na hardwarové peněžence cosignera. Nejvýraznějším
ústupkem oproti desktopovému řešení je delegace komunikace s~Trezorem na
externí aplikaci Trezor Suite Mobile, kterou jsme zvolili jako pragmatický
kompromis mezi rozsahem bakalářské práce a~kvalitou uživatelského prožitku;
důvody i~zvažované alternativy popisujeme v~sekci~\ref{sec:analyza-trezor}.

\section*{Hlavní přínosy práce}

Konkrétní přínosy práce shrnujeme do následujících bodů:

\begin{itemize}
\item \textbf{Mobilní peněženka kombinující multisig, coin control
  a~přímou integraci s~Trezorem bez privátních klíčů na~serveru.}
  Srovnání sedmi existujících řešení v~sekci~\ref{sec:comparison}
  ukázalo, že žádné z~analyzovaných mobilních řešení tuto kombinaci
  nenabízí; nejbližší konkurent BlueWallet~\cite{BlueWallet} podporuje
  multisig na~mobilu, ale komunikace s~hardwarovou peněženkou probíhá
  pouze přes manuální přenos PSBT souborů, bez přímé integrace přes
  deeplink protokol Trezor Connect~\cite{TrezorConnect}. Bezpečnostní
  model jsme ověřili kompletní 2-of-3 multisig transakcí na Bitcoin
  testnetu (viz~kapitolu~\ref{chap:testovani}).

\item \textbf{Engineering deeplink workflow s~Trezor Suite Mobile.}
  Asynchronní podepisování přes Android intenty vyžaduje řešení
  několika netriviálních problémů, jejichž postup není v~dokumentaci
  Trezor Connect~\cite{TrezorConnect} popsán: korelace odpovědí
  s~pending PSBT přes \texttt{requestId}, detekce odpovědi z~jiného
  Trezoru než toho přihlášeného (sentinel \texttt{WRONG\_DEVICE}
  s~auto-logoutem), rozlišení passphrase identity stejného fyzického
  zařízení podle fingerprintu xpub a~obnova stavu po~návratu uživatele
  z~aplikace třetí strany. Konkrétní implementace je popsána
  v~sekci~\ref{subsec:deeplink-workflow}.

\item \textbf{Identifikace implementačních úskalí PSBT a~Trezor
  workflow.} Při ladění jsme narazili na chyby, jejichž diagnostika
  vyžadovala detailní studium specifikací i~bezpečnostních patchů
  firmware Trezoru. Pole
  \texttt{PSBT\allowbreak \_\allowbreak IN\allowbreak \_\allowbreak NON\allowbreak \_\allowbreak WITNESS\allowbreak \_\allowbreak UTXO}
  je vyžadováno i~pro nativní SegWit vstupy počínaje firmware Trezor
  Model T 2.3.1 (resp. Trezor One 1.9.1) z~června 2020 jako mitigace
  BIP-143 SegWit fee útoku zveřejněného Saleemem Rashidem v~březnu
  2020~\cite{TrezorSegwitFix2020}; současný firmware modelů Safe 3, 5
  a~7 toto pravidlo přebírá. Standard BIP-174~\cite{BIP174} přitom
  toto pole pro SegWit vstupy označuje pouze za~volitelné
  (\enquote{\emph{can be present}}). U~deskriptoru
  \texttt{sortedmulti(\ldots)} je třeba podpisy ve~witness stacku
  řadit podle pořadí \emph{odvozených} klíčů, jak explicitně předepisuje
  BIP-383~\cite{BIP383} (\enquote{\emph{after all extended keys are
  derived}}); naivní seřazení xpubs jednou při parsování deskriptoru
  vede k~nevalidní transakci. Konečně předcházení reuse change adres
  vyžádalo vlastní mechanismus sledování stavu napříč rozpracovanými
  PSBT, který není pokryt ani Trezor Connect, ani PSBT formátem.
  Detaily nálezů a~jejich oprav jsou popsány v~sekci~\ref{sec:test-chyby}
  a~mohou posloužit dalším implementátorům.

\item \textbf{Interoperabilita s~peněženkami podporujícími výstupní
  deskriptory bez~importu klíčového materiálu.} Parser deskriptorů
  ve~formátu \texttt{wsh(sortedmulti(\ldots))} podle BIP-380~\cite{BIP380}
  a~BIP-383~\cite{BIP383} přijímá vstup od~jakékoli peněženky, která
  tento otevřený standard respektuje (například Sparrow Wallet, Bitcoin
  Core nebo Specter Desktop). Postup je univerzální: na~desktopu se
  vytvoří multisig konfigurace, exportuje deskriptor a~v~mobilní
  aplikaci se~importuje, čímž obě peněženky pracují nad~totožnou
  množinou adres bez~potřeby přenášet privátní klíče ani seedy mezi
  platformami. Funkčnost jsme ověřili proti Sparrow
  Walletu~\cite{SparrowWallet} jako referenční desktop implementaci
  (viz~sekci~\ref{sec:test-prostredi}).
\end{itemize}

\section*{Možnosti dalšího vývoje}

Aplikace je navržena s~ohledem na další rozšiřování a~několik konkrétních
směrů budoucího vývoje považujeme za přirozené pokračování této práce.

Nejvýznamnějším z~nich je \textbf{přímá komunikace s~hardwarovou peněženkou}
přes Android USB Host API. Současný stav, kdy aplikace deleguje podepisování
na Trezor Suite Mobile, je sice funkční, ale zhoršuje uživatelský komfort
nutností přepínat mezi dvěma aplikacemi a~zavádí závislost na~softwaru třetí
strany. Vlastní implementace by znamenala napsat USB transportní vrstvu
nad 64bajtovým paketovým protokolem Trezoru, kompletní sadu protobuf zpráv
podle aktuálního firmware a~stavový automat pro vícekolovou komunikaci typu
\texttt{SignTx -- TxRequest -- TxAck}, který Trezor používá kvůli omezené
paměti zařízení. Tento rozsah implementace odpovídá samostatnému projektu
přesahujícímu rámec bakalářské práce, proto jsme tento směr vědomě nepřevzali,
ale architektura aplikace
je na něj připravena -- veškerá komunikace s~Trezorem prochází přes interní
abstrakci \texttt{TrezorDeeplinkLauncher} popsanou v~sekci~\ref{sec:impl-detaily},
za~kterou se aktuální deeplink implementace dá nahradit přímou USB
komunikací bez zásahu do zbytku systému.

Druhým přirozeným rozšířením je \textbf{podpora dalších hardwarových peněženek},
zejména Ledger a~Coldcard. Architektura backendu je na~konkrétním modelu
Trezoru nezávislá -- pracuje s~PSBT a~výstupními deskriptory, což jsou
otevřené standardy podporované všemi zmíněnými výrobci. Podpora dalších
zařízení by tak vyžadovala práci pouze na straně mobilní aplikace,
především implementaci odpovídajících transportních vrstev.

Třetím směrem je \textbf{příprava na produkční provoz}. Aplikace již
podporuje přepínání mezi testnetem a~mainnetem přímo z~přihlašovací
obrazovky a~backend automaticky komunikuje s~odpovídajícím blockchainovým
API (Mempool.space pro testnet, Blockstream
Esplora~\cite{BlockstreamEsplora} pro mainnet). Před reálným nasazením
s~finančními prostředky by však bylo vhodné provést bezpečnostní audit
a~rozšířit \textbf{automatizovanou testovací sadu}. Jednotkové testy nad
deterministickými částmi (PSBT a~Trezor Connect parametry, audit trail
podpisů, derivace adres, parser deskriptorů, vystavování a~rotace JWT
tokenů, klasifikace transakcí a~retry logika blockchain klienta) jsme
do~práce zahrnuli a~popsali v~sekci~\ref{sec:test-unit}; chybí
end-to-end pokrytí pomocí Trezor emulátoru, které jsme z~důvodů
diskutovaných v~sekci~\ref{sec:test-pristup} odložili na další vývoj.

Mezi další užitečná rozšíření patří podpora náhradních strategií výběru
UTXO (například algoritmus Branch and Bound), možnost připojit aplikaci
k~vlastnímu Bitcoin uzlu místo veřejných API, podpora schématu Taproot
(BIP-86) a~rozšíření o~zobrazení podrobných statistik a~grafů pohybů
prostředků.
